\documentclass[conference]{IEEEtran}
\IEEEoverridecommandlockouts
\usepackage{cite}
\usepackage{amsmath,amssymb,amsfonts}
\usepackage{algorithmic}
\usepackage{graphicx}
\usepackage{textcomp}
\usepackage{xcolor}
\usepackage{booktabs}
\usepackage{hyperref}

\def\BibTeX{{\rm B\kern-.05em{\sc i\kern-.025em b}\kern-.08em
    T\kern-.1667em\lower.7ex\head{E}\kern-.125emX}}

\begin{document}

\title{Empirical Benchmarking of Classical Machine Learning and Causal Windowing for Multi-Horizon Highway Speed Forecasting on the METR-LA Network}

\author{\IEEEauthorblockN{Member 1 (Lead: Data \& Provenance)}
\IEEEauthorblockA{\textit{Dept. of Computer Science \& AI} \\
\textit{College of Engineering}\\
Email: member1@institution.edu}
\and
\IEEEauthorblockN{Member 2 (Lead: Features \& Baselines)}
\IEEEauthorblockA{\textit{Dept. of Computer Science \& AI} \\
\textit{College of Engineering}\\
Email: member2@institution.edu}
\and
\IEEEauthorblockN{Member 3 (Lead: Models \& Evaluation)}
\IEEEauthorblockA{\textit{Dept. of Computer Science \& AI} \\
\textit{College of Engineering}\\
Email: member3@institution.edu}
\and
\IEEEauthorblockN{Member 4 (Lead: Systems \& Integration)}
\IEEEauthorblockA{\textit{Dept. of Computer Science \& AI} \\
\textit{College of Engineering}\\
Email: member4@institution.edu}
}

\maketitle

\begin{abstract}
Accurate multi-horizon highway traffic speed forecasting is crucial for intelligent transportation systems, congestion mitigation, and commuter routing. While recent research has focused heavily on deep graph neural networks, robust classical machine learning benchmarks and rigorous causal feature validation remain foundational. In this Phase 1 investigation, we present an academically truthful, reproducible classical machine learning framework evaluated on the benchmark METR-LA dataset (207 loop detector stations across Los Angeles County highways). We formalize a strict chronological 70/10/20 train/validation/test partitioning scheme, an explicit missingness protocol distinguishing zero-speed readings from sensor dropouts, and a causal history window builder ($t-11 \dots t$) targeting multi-step forecasting horizons at +15, +30, and +60 minutes. We benchmark Historical Persistence, $L_2$-regularized Ridge Regression, and Gradient Boosted Decision Trees (XGBoost) across 90,149 held-out test target instances. XGBoost achieves a Mean Absolute Error (MAE) of 2.59 mph, 3.17 mph, and 4.01 mph at +15m, +30m, and +60m horizons, respectively, outperforming Persistence (3.00, 3.62, 4.73 mph). Furthermore, a topological feature ablation reveals that incorporating 1st-order road network neighbor lags reduces 60-minute test MAE from 4.02 mph to 3.81 mph. This empirically substantiates the spatial propagation of congestion shockwaves and establishes the formal benchmark for Phase 2 Graph Attention Network (GAT) research.
\end{abstract}

\begin{IEEEkeywords}
Traffic speed forecasting, METR-LA benchmark, causal feature windowing, classical machine learning, XGBoost, multi-horizon evaluation, spatial ablation.
\end{IEEEkeywords}

\section{Introduction}
Intelligent Transportation Systems (ITS) depend heavily on accurate short- and medium-term forecasts of highway traffic speeds. Freeway corridors experience complex spatio-temporal dynamics characterized by sudden bottlenecks, non-linear congestion shockwaves, and diurnal cyclic patterns. 

Under the course mandate \textit{``Traffic Flow Prediction and Signal Control via Graph Attention Networks and Flow Models,''} this paper formalizes the rigorous **Phase 1 Classical Machine Learning and Benchmarking System**. In adherence to strict academic truthfulness, we enforce key scientific boundaries:
\begin{enumerate}
    \item \textbf{Measurement Boundary}: The benchmark METR-LA dataset records loop detector highway speeds in miles per hour (mph). It does not contain vehicular flow/volume or traffic signal phase states. Therefore, no ungrounded claims of live signal actuation or unmeasured flow prediction are asserted.
    \item \textbf{Causal Windowing}: All feature representations are extracted strictly from historical observations ($t-11 \dots t$, representing 60 minutes of elapsed history) to predict forward horizon steps $h \in \{3, 6, 12\}$ (+15m, +30m, +60m).
    \item \textbf{Chronological Partitioning}: A strict chronological 70\% Train, 10\% Validation, 20\% Test temporal split is enforced without target overlap or shuffle leakage.
\end{enumerate}

\section{Related Work}
Classical traffic forecasting literature initially relied on time-series formulations such as Historical Average (HA) and Autoregressive Integrated Moving Average (ARIMA) models \cite{williams2003modeling}. While effective under stationary free-flow regimes, linear models struggle with non-linear transitions during rush-hour congestion shocks. 

The advent of machine learning introduced regularized linear regression (Ridge) \cite{hoerl1970ridge} and ensemble tree-based models such as XGBoost \cite{chen2016xgboost}. In parallel, deep learning research introduced spatial-temporal graph models, notably DCRNN \cite{li2018diffusion} and STGCN \cite{yu2018spatio}, which leverage highway topology matrices. Veli{\v{c}}kovi{\'c} et al. \cite{velickovic2018graph} proposed Graph Attention Networks (GAT), enabling dynamic spatial weighting. Our Phase 1 classical system rigorously establishes the empirical baseline against which Phase 2 GAT architectures will be evaluated.

\section{Dataset Provenance \& Missingness Audit}
\subsection{Data Lineage}
We utilize the canonical METR-LA benchmark introduced in DCRNN \cite{li2018diffusion}, originating from the Caltrans Performance Measurement System (PeMS) \cite{caltrans2012pems}. The raw dataset covers 207 loop detector stations on Los Angeles County freeways over four continuous months (March 1, 2012 to June 27, 2012). Speeds are aggregated into 5-minute intervals, yielding $T = 34,272$ timestamps and $7,094,304$ total speed readings. Cryptographic verification confirms the raw HDF5 file integrity with SHA-256 hash \texttt{64784b...c05}.

\subsection{Missingness \& Validity Accounting}
A comprehensive audit of the 7,094,304 matrix entries reveals:
\begin{itemize}
    \item \textbf{NaN values}: Exactly 0 (null values are represented as zero in the benchmark).
    \item \textbf{Zero readings ($0.0$ mph)}: 575,302 cells (8.11\% of total observations). Loop detector zeros indicate either severe stop-and-go congestion, sensor hardware dropouts, or temporary lane closures.
    \item \textbf{Physiologically impossible speeds ($>100$ mph)}: Exactly 0.
    \item \textbf{Valid speeds ($>0.0$ mph)}: 6,519,002 cells (91.89\%), with an empirical mean of 58.47 mph and standard deviation of 10.30 mph.
\end{itemize}
\textbf{Evaluation Label Policy}: To prevent distortion and undefined divisions during percentage error evaluation, evaluation masks strictly enforce ground-truth target validity ($y_t > 0.0$ mph). For historical feature lags, causal forward-fill imputation ensures complete feature matrices without future information contamination.

\section{Methodology \& Benchmark Formulation}

\subsection{Causal Feature Representation}
For any forecast origin timestamp $t$, the feature representation vector $\mathbf{x}_t \in \mathbb{R}^{D}$ is assembled from:
\begin{enumerate}
    \item \textbf{Speed Lags}: The past 12 five-minute observations: $[y_{t-11}, y_{t-10}, \dots, y_t]$.
    \item \textbf{Rolling Window Statistics}: Historical mean, standard deviation, minimum, maximum, and linear momentum trend ($\bar{y}_{t-2:t} - \bar{y}_{t-11:t-9}$).
    \item \textbf{Cyclical Calendar Encodings}: Continuous diurnal and weekly sinusoidal projections:
    \begin{equation}
        \sin\left(\frac{2\pi \cdot \text{TOD}}{1440}\right), \cos\left(\frac{2\pi \cdot \text{TOD}}{1440}\right), \sin\left(\frac{2\pi \cdot \text{DOW}}{7}\right), \cos\left(\frac{2\pi \cdot \text{DOW}}{7}\right)
    \end{equation}
    \item \textbf{Topological Spatial Lags (Ablation)}: Recent 3 historical speed observations from the top-$K$ ($K=3$) nearest 1st-order road network graph neighbors derived from the thresholded Gaussian distance matrix $\mathbf{A} \in \mathbb{R}^{207 \times 207}$.
\end{enumerate}

\subsection{Partitioning \& Preprocessing Integrity}
Partitions are chronologically partitioned without shuffling:
\begin{itemize}
    \item \textbf{Train (70\%)}: 23,990 steps (2012-03-01 to 2012-05-23).
    \item \textbf{Validation (10\%)}: 3,427 steps (2012-05-23 to 2012-06-04).
    \item \textbf{Test (20\%)}: 6,855 steps (2012-06-04 to 2012-06-27).
\end{itemize}
All feature scaling transformations ($\mathbf{z} = (\mathbf{x} - \boldsymbol{\mu}_{\text{train}}) / \boldsymbol{\sigma}_{\text{train}}$) are fitted exclusively on the training partition and applied out-of-sample to validation and test partitions.

\subsection{Benchmark Models}
\textbf{1. Historical Persistence}:
\begin{equation}
    \hat{y}_{t+h} = y_t, \quad \forall h \in \{3, 6, 12\}
\end{equation}
\textbf{2. Multi-Horizon Ridge Regression}: Direct regularized linear models minimizing $L_2$ penalized loss:
\begin{equation}
    \mathcal{L}_{\text{Ridge}}(\mathbf{w}_h) = \sum_{i=1}^{N_{\text{train}}} (y_{i,t+h} - \mathbf{w}_h^T \mathbf{x}_{i,t})^2 + \alpha_h \|\mathbf{w}_h\|_2^2
\end{equation}
where $\alpha_h$ is tuned on validation MAE ($\alpha_{15m}=100.0, \alpha_{30m}=0.1, \alpha_{60m}=0.1$).

\textbf{3. Multi-Horizon XGBoost}: Gradient boosted decision tree ensembles minimizing regularized squared error via second-order Taylor expansion:
\begin{equation}
    \mathcal{L}_{\text{XGB}}^{(m)} = \sum_{i} \left[ g_i f_m(\mathbf{x}_i) + \frac{1}{2} h_i f_m^2(\mathbf{x}_i) \right] + \Omega(f_m)
\end{equation}
Hyperparameters: $n_{\text{estimators}}=250$, learning rate $\eta=0.05$, max depth $d=6$, subsample ratio $0.8$, with validation early stopping.

\subsection{Evaluation Metrics}
All models are evaluated on identical valid target masks across three horizons ($h=3, 6, 12$, corresponding to 15, 30, and 60 minutes ahead):
\begin{align}
    \text{MAE} &= \frac{1}{|\mathcal{V}|} \sum_{i \in \mathcal{V}} |y_i - \hat{y}_i| \\
    \text{RMSE} &= \sqrt{\frac{1}{|\mathcal{V}|} \sum_{i \in \mathcal{V}} (y_i - \hat{y}_i)^2} \\
    \text{MAPE} &= \frac{1}{|\mathcal{V}|} \sum_{i \in \mathcal{V}} \frac{|y_i - \hat{y}_i|}{\max(y_i, \epsilon)} \times 100\%
\end{align}
with documented $\epsilon = 2.0$ mph to prevent numerical instability during near-zero traffic jams.

\section{Experimental Results}

\begin{table}[t]
\caption{Held-Out Test Set Performance Across Horizons (Pilot Cohort, 90,149 Valid Target Instances)}
\label{tab:metrics_horizon}
\centering
\begin{tabular}{llcccc}
\toprule
\textbf{Model} & \textbf{Horizon} & \textbf{MAE (mph)} & \textbf{RMSE (mph)} & \textbf{MAPE (\%)} \\
\midrule
Persistence & 15 min & 3.0011 & 5.5959 & 7.12\% \\
Persistence & 30 min & 3.6157 & 7.2034 & 9.04\% \\
Persistence & 60 min & 4.7302 & 9.6660 & 12.68\% \\
\midrule
Ridge       & 15 min & 2.8789 & 5.3466 & 7.45\% \\
Ridge       & 30 min & 3.6050 & 6.8440 & 10.12\% \\
Ridge       & 60 min & 4.9066 & 8.9476 & 15.02\% \\
\midrule
\textbf{XGBoost} & \textbf{15 min} & \textbf{2.5909} & \textbf{4.8718} & \textbf{6.68\%} \\
\textbf{XGBoost} & \textbf{30 min} & \textbf{3.1668} & \textbf{6.0589} & \textbf{8.69\%} \\
\textbf{XGBoost} & \textbf{60 min} & \textbf{4.0083} & \textbf{7.5269} & \textbf{11.64\%} \\
\bottomrule
\end{tabular}
\end{table}

\subsection{Multi-Horizon Performance}
Table \ref{tab:metrics_horizon} summarizes the empirical test results. XGBoost achieves superior accuracy across all evaluated horizons. At 15 minutes ahead, XGBoost achieves an MAE of 2.59 mph, representing a 13.7\% relative error reduction over the Persistence benchmark (3.00 mph). At 60 minutes ahead, error increases across all models as the forecast horizon outpaces short-term autocorrelation; however, XGBoost maintains an MAE of 4.01 mph compared to 4.73 mph for Persistence and 4.91 mph for Ridge.

\subsection{Peak vs. Off-Peak Stratified Error Analysis}
Table \ref{tab:metrics_peak} demonstrates the dramatic impact of rush-hour congestion dynamics. During morning (07:00-09:30) and evening (16:00-19:00) weekday peak periods, error inflates significantly for all baselines due to rapid non-linear transitions between free-flow and stop-and-go states. Ridge regression exhibits severe peak degradation at 60 minutes (MAE 6.47 mph), whereas XGBoost exhibits greater resilience (MAE 5.09 mph).

\begin{table}[t]
\caption{Stratified Peak vs. Off-Peak Test Error (XGBoost vs Baselines)}
\label{tab:metrics_peak}
\centering
\begin{tabular}{llcccc}
\toprule
\textbf{Model} & \textbf{Horizon} & \textbf{Peak MAE} & \textbf{Off-Peak MAE} & \textbf{Peak RMSE} \\
\midrule
Persistence & 15 min & 4.2384 & 2.6517 & 7.4206 \\
Persistence & 60 min & 6.4022 & 4.2573 & 12.0620 \\
\midrule
Ridge       & 15 min & 4.0851 & 2.5379 & 6.9458 \\
Ridge       & 60 min & 6.4729 & 4.4637 & 11.2387 \\
\midrule
\textbf{XGBoost} & \textbf{15 min} & \textbf{3.6596} & \textbf{2.2887} & \textbf{6.2559} \\
\textbf{XGBoost} & \textbf{60 min} & \textbf{5.0863} & \textbf{3.7034} & \textbf{9.2136} \\
\bottomrule
\end{tabular}
\end{table}

\section{Spatial Feature Ablation \& Phase 2 Outlook}
To rigorously evaluate the utility of the road graph topology prior to full deep graph modeling, we executed an ablation study comparing isolated single-sensor historical lags (21 features) against models augmented with recent historical lags from the top-3 nearest 1st-order graph neighbors (30 features). 

As shown in Table \ref{tab:ablation}, adding 1st-order topological neighbors yields consistent error reductions across all horizons, with the most pronounced improvement at +60 minutes where MAE decreases from 4.02 mph to 3.81 mph (a 5.2\% relative improvement) and RMSE drops from 7.54 mph to 7.03 mph.

\begin{table}[t]
\caption{Road Network Spatial Lag Feature Ablation (XGBoost Regressor)}
\label{tab:ablation}
\centering
\begin{tabular}{lcccc}
\toprule
\textbf{Configuration} & \textbf{Horizon} & \textbf{MAE (mph)} & \textbf{RMSE (mph)} & \textbf{MAPE (\%)} \\
\midrule
Isolated Sensor Lags   & 15 min & 2.5949 & 4.8789 & 6.69\% \\
\textbf{With Spatial Neighbors} & \textbf{15 min} & \textbf{2.5464} & \textbf{4.6589} & \textbf{6.46\%} \\
\midrule
Isolated Sensor Lags   & 30 min & 3.1736 & 6.0692 & 8.71\% \\
\textbf{With Spatial Neighbors} & \textbf{30 min} & \textbf{3.0727} & \textbf{5.7577} & \textbf{8.32\%} \\
\midrule
Isolated Sensor Lags   & 60 min & 4.0211 & 7.5400 & 11.70\% \\
\textbf{With Spatial Neighbors} & \textbf{60 min} & \textbf{3.8112} & \textbf{7.0303} & \textbf{10.96\%} \\
\bottomrule
\end{tabular}
\end{table}

\subsection{Architectural Blueprint for Phase 2 GAT}
While classical feature concatenation demonstrates the value of spatial context, it imposes fixed, static edge weights. In Phase 2, we will implement dynamic Spatio-Temporal Graph Attention Networks (ST-GAT). The attention coefficient $\alpha_{ij}^{(t)}$ between sensors $i$ and $j$ will be computed dynamically as:
\begin{equation}
    \alpha_{ij}^{(t)} = \frac{\exp\left(\text{LeakyReLU}\left(\mathbf{a}^T [\mathbf{W}h_i^{(t)} \parallel \mathbf{W}h_j^{(t)}]\right)\right)}{\sum_{k \in \mathcal{N}_i} \exp\left(\text{LeakyReLU}\left(\mathbf{a}^T [\mathbf{W}h_i^{(t)} \parallel \mathbf{W}h_k^{(t)}]\right)\right)}
\end{equation}
enabling the network to dynamically adapt its spatial reception field according to prevailing traffic shockwave velocities.

\section{Conclusion}
This study presented a fully reproducible, academically truthful Phase 1 classical machine learning system for METR-LA highway traffic speed forecasting. By enforcing strict chronological boundaries, an explicit missingness protocol, and causal window construction, we established genuine empirical benchmarks for Historical Persistence, Ridge Regression, and XGBoost. The superior performance of XGBoost (2.59 mph at 15m) and the empirical validation of topological neighbor lags (3.81 mph at 60m) provide the rigorous scientific foundation required for Phase 2 Graph Attention Network development.

\bibliographystyle{IEEEtran}
\bibliography{references}

\end{document}
